\section{Commensurability, named lattices, and endomorphisms}
\label{sec:examples}

The category calculation gives concise consequences with different
mathematical content. Untwisted correspondence existence detects abstract
commensurability. Stable isomorphism recovers an entire group and a
specified map. A finite-index endomorphism requires one additional
classical covolume argument; its two module dimensions are not initially
reciprocal.

\subsection{Untwisted correspondences and finite-type neighbors}

\begin{corollary}\label{cor:commensurability}
Let $\Gamma\in\mathscr K$ and let $\Lambda$ be any countably infinite
ICC group. There is a nonzero bifinite $L(\Gamma)$--$L(\Lambda)$
correspondence if and only if $\Gamma$ and $\Lambda$ are abstractly
commensurable.
\end{corollary}
\begin{proof}
A nonzero correspondence has a nonzero full-germ orbit in
Theorem~\ref{thm:correspondences}. Its stabilizer is the graph of an
actual isomorphism between finite-index subgroups. Conversely such an
isomorphism, with its one-dimensional trivial fibre representation,
gives the model in Equation~\eqref{eq:intro-model}.
\end{proof}

With the finite-type convention of Section~\ref{sec:introduction},
the endpoint theorem specializes as follows.
\begin{corollary}\label{cor:untwisted-neighbors}
For $\Gamma\in\mathscr K$, a separable $\mathrm{II}_1$ factor $P$ is
joined to $L(\Gamma)$ by a nonzero bifinite correspondence if and only if
$P\cong L_\omega(\Lambda)^t$ for $t>0$, an ICC group $\Lambda$
abstractly commensurable with $\Gamma$, and a finite-type cocycle
$\omega$ on $\Lambda$.
\end{corollary}
\begin{proof}
Theorem~\ref{thm:endpoints} and the final projective-induction argument
in its proof give necessity. For sufficiency, restrict a projective
representation of multiplier $\omega$ to the matched subgroup and
conjugate it; its multiplier is exactly $1/\delta^*\omega$. The graph
model and the amplification equivalence give the required correspondence.
\end{proof}

For arbitrary source twist this specialization must be replaced by the
ratio in Theorem~\ref{thm:endpoints}. In particular, its direct proof
admits the identity correspondence of $L_\mu(\Gamma)$ whatever the
individual cocycle $\mu$ may be.

\subsection{Examples and specified canonical maps}

The familiar arithmetic examples are
$\mathrm{SL}_3(\mathbb Z)$, its principal congruence subgroup
\[
 \Gamma(3)=\ker\bigl(\mathrm{SL}_3(\mathbb Z)
                    \longrightarrow\mathrm{SL}_3(\mathbb Z/3\mathbb Z)\bigr),
\]
all their finite-index subgroups, $\mathrm{PSL}_n(\mathbb Z)$ for
$n\geq4$, and $\mathrm{PSp}_4(\mathbb Z)$, with the latter two
symbols denoting the quotient by the center of the integral matrix
group. The arithmetic lattice theorem places these projective groups in
the corresponding center-free real groups \cite{BorelHarishChandra1962}; the finite-covolume statement appears also in
\cite[Theorem~1(iv)]{BorelHarishChandra1961}.
For even $n$, $\mathrm{PSL}_n(\mathbb R)$ is the identity root-generated
subgroup of $\mathrm{PGL}_n(\mathbb R)$; for odd $n$ it is all of it.
Likewise $\mathrm{PSp}_4(\mathbb R)$ is the root-generated adjoint
symplectic group. These ambient groups have real rank at least two and
property~$(T)$ \cite{BHV2008}.

These projective lattices are ICC. A finite conjugacy class gives a
finite-index centralizer. That subgroup is Zariski dense by Borel density,
so the element centralizes the connected ambient algebraic group and is
trivial in its adjoint form \cite{Borel1960}. Equivalently, the centralizer formulation in
\cite[Corollary~3]{Furstenberg1976} applies directly to the finite-index
lattice and gives the ambient center. For
$\mathrm{SL}_3(\mathbb Z)$ the real center is already trivial.
A finite-index subgroup of an ICC group remains ICC: its finite
conjugacy class would yield a finite class in the ambient group after
using finitely many coset representatives. Thus every group just listed
belongs to $\mathscr K$.

For any such group $\Gamma$, any countably infinite ICC group $\Lambda$,
and any specified isomorphism
\[
 \theta:L(\Gamma)\longrightarrow pM_k(L(\Lambda))p,
\]
Theorem~\ref{thm:stable} yields
\[
 (\Tr_k\otimes\tau)(p)=1,\qquad
 \theta(u_g)=\chi(g)Wv_{\delta(g)}W^*,\qquad
 W^*W=1,\quad WW^*=p,
\]
where $W\in M_{k,1}(L(\Lambda))$,
$\delta:\Gamma\to\Lambda$ is an actual group isomorphism, and
$\chi:\Gamma\to\mathbb T$ is a character. This includes the
specified column in the specified corner. The character becomes the
cochain of Equation~\eqref{eq:intro-cochain} when scalar twists are
present. In particular all the corresponding twisted factors have
fundamental group $\{1\}$.

Definition~\ref{def:class} also includes ICC lattices in the permitted
finite-place groups, mixed real and finite-place products, and products
with quaternionic or Cayley real factors of property~$(T)$. Finite
products of the preceding discrete examples give reducible cases:
a nonidentity element has a nonidentity coordinate with infinite
conjugacy class. More generally the companion exhaustion theorem already
allows arbitrary reducible lattices, nonuniform lattices, and torsion.
No repetition of the individual geometric proofs is required for these
examples.

\begin{corollary}[Outer automorphisms]\label{cor:outer-automorphisms}
For $\Gamma\in\mathscr K$, there is an isomorphism of abstract groups
\[
 \operatorname{Out}(L(\Gamma))
 \cong\Hom(\Gamma,\mathbb T)\rtimes\operatorname{Out}(\Gamma),
 \qquad [\delta]\cdot\chi=\chi\circ\delta^{-1},
\]
where $\operatorname{Out}$ denotes automorphisms modulo inner
automorphisms.
\end{corollary}
\begin{proof}
Put $T_\chi(u_g)=\chi(g)u_g$ and $U_\delta(u_g)=u_{\delta(g)}$.
These maps extend normally by their trace-preserving basis unitaries.
Characters are invariant under inner group automorphisms, and
$U_\delta T_\psi U_\delta^{-1}=T_{\psi\circ\delta^{-1}}$.
Thus $(\chi,[\delta])\mapsto[T_\chi U_\delta]$ is a well-defined
homomorphism for the displayed left action. It is surjective by the
untwisted unit-scale case of Theorem~\ref{thm:stable}, whose cochain
is then a character.

For injectivity, suppose $T_\chi U_\delta=\Ad(w)$, and write
$b=\chi\circ\delta$ and $w=\sum_x a_xu_x$ in $L^2(L(\Gamma))$.
The relation
\[
 u_{\delta(g)}wu_g^{-1}=b(g)^{-1}w
\]
shows that $|a_x|$ is constant on each orbit of
$x\mapsto\delta(g)xg^{-1}$. A nonzero coefficient has a finite orbit
by square summability. Its stabilizer $H$ has finite index, and
$\delta(h)=xhx^{-1}$ for all $h\in H$.
Set $\beta=\Ad(x^{-1})\circ\delta$ and let $N$ be the normal core
of $H$. Since $\beta$ fixes $N$ pointwise,
$g^{-1}\beta(g)$ centralizes $N$ for every $g\in\Gamma$.
The centralizer of a finite-index subgroup is trivial by ICC, as in
Lemma~\ref{lem:calculus-germ-test}; hence $\delta=\Ad(x)$.
Now $v=u_x^*w$ implements the pure character twist $b$.
Its Fourier-coefficient absolute values are constant on conjugacy
classes, so ICC makes $v$ scalar. Therefore $b=1$, whence $\chi=1$,
and $[\delta]=1$, proving injectivity.
\end{proof}

\subsection{Recovery from the reduced group C*-algebra}

Write $C_r^*(\Gamma)$ for the operator-norm closure of the group algebra
in the left regular representation.
\begin{corollary}[Reduced C*-algebra recovery]\label{cor:reduced-cstar}
Let $\Gamma\in\mathscr K$ and suppose that $C_r^*(\Gamma)$ has a
unique tracial state. If $\Lambda$ is any discrete group and
there is a complex-linear $*$-isomorphism
$C_r^*(\Gamma)\cong C_r^*(\Lambda)$, then
$\Gamma\cong\Lambda$.
\end{corollary}
\begin{proof}
Let $\alpha:C_r^*(\Gamma)\to C_r^*(\Lambda)$ be the given
$*$-isomorphism, and denote the canonical traces by $\tau_\Gamma$
and $\tau_\Lambda$. The isomorphism is unital, and uniqueness gives
$\tau_\Lambda\circ\alpha=\tau_\Gamma$. Consequently $\alpha$
extends isometrically to a unitary between the trace-GNS Hilbert
spaces. This unitary intertwines left multiplication, so conjugation by
it extends $\alpha$ to a normal isomorphism
$L(\Gamma)\cong L(\Lambda)$.

The source reduced algebra is separable. Distinct canonical group
unitaries in the target have trace $2$-norm distance $\sqrt2$, and hence
operator-norm distance at least $\sqrt2$; thus $\Lambda$ is countable.
Moreover, $L(\Lambda)$ is an infinite-dimensional factor, so $\Lambda$
is infinite and ICC. Indeed, a nonidentity finite conjugacy class would
give a nonscalar central class sum. Theorem~\ref{thm:stable}, with both
cocycles trivial and both amplification scales equal to one, now gives
$\Gamma\cong\Lambda$.
\end{proof}

In particular this applies to $\Gamma=\mathrm{SL}_3(\mathbb Z)$:
its reduced group C*-algebra has a unique tracial state by
\cite[discussion following Corollary~5]{Bekka2007}.

\subsection{Finite-index endomorphisms of the special linear factor}

We state the precise classical fact needed in the final deduction.
The finite-index image hypothesis is part of it.
\begin{lemma}\label{lem:sl3-covolume}
An injective homomorphism
$\alpha:\mathrm{SL}_3(\mathbb Z)\to\mathrm{SL}_3(\mathbb Z)$
whose image has finite index is surjective.
\end{lemma}
\begin{proof}
Put $G=\mathrm{SL}_3(\mathbb R)$ and $\Gamma=\mathrm{SL}_3(\mathbb Z)$.
The group $G$ is connected, center-free, simple, of real rank two, and
has no compact factor. Both $\Gamma$ and $\alpha(\Gamma)$ are
irreducible lattices in $G$. The classical strong-rigidity theorem
extends their isomorphism to an analytic automorphism $\Phi$ of $G$;
see \cite[Theorem~B]{Prasad1973}. This form of strong rigidity allows
nonuniform lattices and torsion; its excluded real rank-one
$\mathrm{SL}_2(\mathbb R)$ case does not occur here. The theorem
combines the uniform rigidity of Mostow with the nonuniform higher-rank
work of Margulis and Raghunathan and Prasad's extension, as explained
in that reference.

The differential $D=d\Phi_e$ preserves the nondegenerate Killing form
$B$ on the Lie algebra: $D^{\mathsf T}BD=B$. Hence $|\det D|=1$.
A left-invariant Haar density is therefore preserved at the identity
and, by translation, everywhere. Consequently
\[
 \operatorname{covol}(\alpha(\Gamma))
 =\operatorname{covol}(\Phi(\Gamma))
 =\operatorname{covol}(\Gamma).
\]
On the other hand the first covolume equals
$[\Gamma:\alpha(\Gamma)]\operatorname{covol}(\Gamma)$.
The common covolume is positive and finite, so the index equals one.
\end{proof}

\begin{corollary}\label{cor:sl3-endomorphism}
Every normal unital complex-linear $*$-endomorphism $\theta$ of
$M=L(\mathrm{SL}_3(\mathbb Z))$ with $[M:\theta(M)]<\infty$
is surjective.
\end{corollary}
\begin{proof}
The endomorphism is injective because its kernel is an ultraweakly
closed two-sided ideal in a factor, and it preserves the normalized trace. Equip $H=L^2(M)$ with its usual
right action and the left action $x\xi=\theta(x)\xi$. It is bifinite,
with
\[
 \dim H_M=1,\qquad \dim_M H=[M:\theta(M)].
\]
Its actual full-germ decomposition and
Equation~\eqref{eq:calculus-summand-dimensions} show that right dimension
one leaves exactly one orbit, a one-dimensional fibre, and
$A_d=\mathrm{SL}_3(\mathbb Z)$. The graph is therefore an isomorphism
$\delta:\mathrm{SL}_3(\mathbb Z)\to B_d$, where $B_d$ has finite
index in the same group. Lemma~\ref{lem:sl3-covolume} makes $B_d$ the
whole group. The left dimension formula now gives $[M:\theta(M)]=1$,
so $\theta(M)=M$.
\end{proof}
